\chapter{Аналитический раздел}

\section{Аффективная робототехника}

Развитие роботизированных технологий влечёт за собой неизбежное возникновение вопросов, связанных с взаимодействием человека и робота, а также с восприятием робота как полноценного участника этого взаимодействия~\cite{akmaev2025dynamics}.
Исследования в области ЧМВ показывают, что люди при взаимодействии с роботами используют те же социальные и коммуникативные механизмы, что и при взаимодействии друг с другом~\cite{park2012lawofattraction}.

Аффективная робототехника --- сфера деятельности, которая является пересечением между областями аффективных вычислений и человеко-машинной коммуникации~\cite{spitale2022affective}. % вывод из имеющейся информации
Согласно с ограничениями применимости ДС в данной работе и классификацией сервисных роботов по их сферам применимости, созданной международной федерацией робототехники~\cite{international_federation_of_robotics}, можно выделить следующие области применения аффективных ДС:

\begin{itemize}
	\item роботы-компаньоны (AC21) --- взаимодействие с людьми и их развлечение;
	\item образование (AC22) --- помощь в обучении детей, моделировании явлений, дистанционном участии в уроках и поддержке учителей~\cite{affect-rob-in-edu};
	\item медицина (AC39, AP63) --- уход за пожилыми людьми и поддержка людей с когнитивными нарушениями~\cite{affect-rob-in-med, affect-rob-in-med-2}.
\end{itemize}

Цель аффективной робототехники --- приблизить взаимодействие человека с машиной к человеческому опыту и восприятию, полученному в межличностном взаимодействии. % определение из паблика в вк (но раньше был и сайт)
Средством достижения этой цели может быть введение такой функциональности роботов, которая позволит им обмениваться с людьми не только информацией, но и аффективной составляющей её восприятия --- аналогом эмоциональной окраски информации человеком или его отношения к ней.
Для формирования такого ЧМВ требуется придать роботам способность распознавать и интерпретировать человеческие эмоции и/или проявлять социальное и эмоциональное поведение.

\section{Робот Ф-2}\label{F2arch}

Робот Ф-2 --- исследовательский проект, в рамках которого изучаются способы, как сделать роботов более привлекательными для человека~\cite{f2_robot}.
Робот способен воспроизводить эмоциональные и рациональные реакции, используя жесты, речь, мимику.
В роботе выделяются три основных компонента (в рамках данной работы не рассматриваются компоненты, связанные с компьютерным зрением и тактильным восприятием)~\cite{kotov_f2_architecture}: лингвистический, сценарный и управляющий.

\subsection{Лингвистический компонент}

Задачей лингвистического компонента (ЛК) является обработка входящего сообщения на естественном языке (текстового или аудио) с целью построения правильной реакции робота на него.
ЛК разделен на 3 модуля, которые работают последовательно~\cite{kotov_f2_architecture}.

\begin{enumerate}
	\item Морфологический модуль --- каждой словоформе в тексте приписываются морфологические и семантические признаки~\cite{zaydelman2018_f2_text}. 
	В основе модуля лежит словарь из 100 тысяч лексем на основе проекта OpenCorpora~\cite{open_coprora}.
	\item Синтаксический модуль --- для каждого предложения строятся одно или несколько деревьев синтаксических зависимостей. 
	Для этого используется словарь формализованных правил русского языка на языке syntXML~\cite{kotov2020_f2_architecture}.
	\item Семантический модуль --- для каждого дерева синтаксических зависимостей строится его семантическое представление: узлам дерева назначаются валентности и семантические признаки, связанные с валентностью. Узлы дерева помещаются в линейную структуру данных на основе постулата, что корнем всегда является глагол в валентности предиката~\cite{kotov2020_f2_architecture}.
\end{enumerate}

Полученное семантическое представление поступает в компонент сценариев.

\subsection{Компонент сценариев}

Поступающее на вход семантическое представление сравнивается с набором сценариев.
Д-сценарии есть доминантные сценарии --- аффективные сценарии реагирования, базовые аналоги эмоций, p-сценарии --- рациональные сценарии реагирования~\cite{kotov_nosovets2023_derived_representation}.
Каждый сценарий включает семантические структуры --- множества признаков, распределённых по валентностям.
Следовательно, для каждой пары вида <семантическое представление, сценарий> вычисляется мера близости с посылкой сценария, также семантическим представлением факта.
Значение меры зависит от числа совпавших семантических признаков в тождественно равных валентностях для пары фактов.
На основе меры близости и чувствительности сценариев вычисляется активизация каждого сценария~\cite{zinina2018_f2_rec}.

Сценарий, получивший при активации наибольший вес, выбирается для реализации.
Для коммуникативной цели ответа --- второй части сценария --- из базы реакций выбирается произвольная реакция из числа сопоставленных данной цели.
Мультимодальная реакция включает жесты, элементы мимики и текст, оформленные в формате BML (Behavior Markup Language).
Кадр поведения передаётся в компонент управления для исполнения на роботе.

\subsection{Компонент управления}

Компонент управления отвечает за выполнение поступивших BML-пакетов с помощью жестов, мимики, речи~\cite{kotov_f2_architecture}.
Речевая составляющая воспроизводится путём синтеза поступающего высказывания c использованием YandexSpeech API.
Жестовая составляющая выполняется на основе перемещения сервомоторов робота по контрольным точкам, указанным в пакете команд, описывающих жест~\cite{Arinkin2016VKR}.

\section{Подходы к описанию персональных черт}

Черта --- это диспозиция (склонность) вести себя определенным образом, которая проявляется в поведении человека в широком спектре ситуаций~\cite{PervinJohn2001Personality}.
Рассматривая подходы к описанию личности человека посредством выделения некоторого набора персональных черт, нужно упомянуть, что такая концепция личности имеет иерархический характер: от самых общих, определяющих глобальные тенденции поведения --- суперфакторов, до более частных и конкретных, проявляющихся в специфических ситуациях.
Иными словами, черты высшего уровня включают множество взаимосвязанных характеристик нижнего уровня, образуя многоуровневую структуру~\cite{PervinJohn2001Personality}.

\subsection{Личностные черты по Г. У. Олпорту}

В выделенных Г. У. Олпортом уровнях черт прослеживается следующая иерархия~\cite{PervinJohn2001Personality}:

\begin{itemize}
	\item кардинальные черты --- доминирующие диспозиции, которые пронизывают всю жизнь человека и проявляются практически во всех его действиях;
	\item центральные черты --- основные характеристики, описывающие поведение человека в повседневной жизни: честность, доброта, общительность, настойчивость;
	\item вторичные черты --- менее устойчивые черты, проявляющиеся в ограниченных ситуациях и отражающие индивидуальные предпочтения, привычки и вкусы.
\end{itemize}

Г. У. Олпорт утверждал, что черта личности позволяет понять, как человек обычно ведёт себя во многих ситуациях, но не предопределяет поведение в каждой отдельной ситуации.
Поэтому для полного понимания поведения индивида необходимо также учитывать понятие ситуации~\cite{PervinJohn2001Personality}.
Несмотря на проведённую теоретическую работу, Олпорт почти не проводил исследований на выявление конкретных черт человека~\cite{PervinJohn2001Personality}.

\subsection{Трёхфакторная теория Г. Ю. Айзенка}

Г. Ю. Айзенк выделяет четыре уровня в структуре личности: уровень конкретных реакций, уровень привычных реакций, уровень черт и уровень суперфакторов~\cite{hjellegger-zigler, PervinJohn2001Personality}.
В результате дважды проведённой процедуры факторного анализа Г. Ю. Айзенк выделил 2 суперфактора личности --- экстраверсию и нейротизм.
Каждый из них характеризует непрерывный диапазон различных степеней выраженности черты: экстраверсия --- интроверсия, нейротизм --- стабильность~\cite{theory_of_personality_eysenck}.

Г. Ю. Айзенк описывает человека с высоким уровнем экстраверсии как общительного, импульсивного, активного, в то время как интроверт рассматривается как человек замкнутый, любящий уединение.
Высокий уровень нейротизма связан с тревогой и беспокойством, низкий --- со спокойствием, уравновешенностью и сдержанностью~\cite{theory_of_personality_eysenck}.
Крайние значения суперфакторов позволяют выделить 4 типа личности, которые соотносятся с 4 типами темперамента, описанных Гиппократом и Галленом (рисунок~\ref{img:eysenck-temperament})~\cite{theory_of_personality_eysenck, PervinJohn2001Personality}.

\includeimage
{eysenck-temperament}
{f}
{h}
{0.65\textwidth}
{Соотношение типов темперамента с суперчертами, выделенными Г. Ю. Айзенком}

Впоследствии к этим двум параметрам личности был добавлен третий --- психотизм, высокий уровень которого характеризует человека как холодного, равнодушного к другим, склонного к риску и противника общепринятых социальных норм~\cite{theory_of_personality_eysenck}.
Айзенк также разработал множество вопросников для измерения личности по этим параметрам: Личностный список модели, Личностный список Айзенка, Личностный вопросник Айзенка~\cite{PervinJohn2001Personality}.

\subsection{Подход Р. Б. Кеттела}

В рамках своей теории Р. Б. Кеттел предложил две классификации личностных черт.
Согласно первой, черты делятся на черты способностей, черты темперамента и динамические черты~\cite{PervinJohn2001Personality}.
Первые характеризуют индивидуальные навыки и умения, вторые отражают эмоциональные особенности и стиль поведения, третьи описывают мотивационную сферу личности.

Согласно второй классификации, черты подразделяются на поверхностные и глубинные~\cite{PervinJohn2001Personality}.
Поверхностные черты являются наблюдаемыми особенностями поведения, которые внешне кажутся взаимосвязанными.
Глубинные черты представляют собой фундаментальные структуры личности, стоящие за наблюдаемыми проявлениями.
Для обнаружения глубинных черт личности Кеттелом были обработаны данные о реальном поведении, данные из опросников и данные поведенческих тестов~\cite{PervinJohn2001Personality, theory_of_personality_cattell}.
В результате многократного применения факторного анализа было выделено 16 факторов~\cite{hjellegger-zigler}, для измерения которых был создан 16-факторный опросник~\cite{cattell_questionnaire}.

\subsection{Большая пятёрка черт}

Пятифакторная модель личности, известная также как Большая пятёрка черт (БПЧ), возникла как обобщение прежде известных диспозициальных теорий личности~\cite{PervinJohn2001Personality}.
Основой <<Большой пятёрки>> является фундаментальная лексическая гипотеза --- гипотеза о том, что со временем наиболее важные индивидуальные различия были закодированы в виде отдельных слов в естественном языке~\cite{big_five_traits}.
Термин <<Большая пятёрка>> был введён Л. Голдбергом в своей работе 1981 года~\cite{goldberg1981language}.

\begin{itemize}
	\item Экстраверсия показывает, насколько человек любит общение, полон энергии и положительных эмоций.
	\item Доброжелательность отражает социально-ориентированную направленность личности, терпимость и готовность учитывать интересы других.
	\item Добросовестность показывает, насколько добросовестно человек выполняет обязанности, насколько он организован и мотивирован.
	\item Нейротизм отражает склонность испытывать негативные эмоции и уровень эмоциональной стабильности.
	\item Открытость опыту отражает интерес человека к новому: идеям, людям, местам, способам действия.
\end{itemize}

Каждый из вышеперечисленных факторов включает в себя 6 граней --- дополнительных, более конкретных аспектов, без которых понимание основных факторов было бы неполным~\cite{big_five_traits}.
Детализированная классификация приведена с обозначениями на рисунке~\ref{img:big-five-traits}.

\includeimage
{big-five-traits}
{f}
{H}
{0.85\textwidth}
{Детализированная классификация черт <<Большой пятёрки>>}

Существует множество способов измерения Большой пятёрки.
Наиболее популярным является коммерческий опросник NEO PI-R, состоящий из 240 вопросов.
Существуют также свободно доступные аналоги этого теста, такие как IPIP-NEO-300 и IPIP-NEO-120.
Один из последних методов измерения большой пятёрки --- BFI-2, разработанный Оливером Джоном и Кристофером Сото в 2015 году~\cite{big_five_inventory}.

\subsection{Типология Майерс-Бриггс}

Типология Майерс-Бриггс (Myers---Briggs Type Indicator, MBTI) --- это психологическая типология личности, разработанная Изабель Бриггс Майерс и Кэтрин Бриггс как дополнение типологии Карла Юнга~\cite{hjellegger-zigler}.
MBTI широко используется в сфере личностного развития, управления персоналом, образования и командного взаимодействия.
Главная идея MBTI заключается в том, что любой вид человеческой личности можно описать, используя один из 16 типов.
В основе этого лежат 4 характеристики, каждая из которых имеет 2 противоположности~\cite{Pittenger1993MeasuringMBTI}: экстраверсия --- интроверсия, сенсорика --- интуиция, мышление --- чувство, суждение --- восприятие.

Комбинация четырёх характеристик образует 16 типов личности.
Каждый тип представляет способ восприятия, мышления и взаимодействия с миром, однако типологический характер MBTI ограничивает точность задания промежуточных состояний личности.

\section{Сравнение подходов к описанию персональных черт}

Сравнение ПОПЧ проведено по 5 критериям, результаты представлены в таблице~\ref{tbl:exist_sol_transposed}.
В качестве критериев сравнения ПОПЧ были взяты следующие:

\begin{itemize}
	\item \textbf{K1} --- опубликованы свидетельства о применимости данного подхода при моделировании различного поведения персонажей ДС или роботов-собеседников;
	\item \textbf{K2} --- вариативность моделируемых личностей;
	\item \textbf{K3} --- наличие теста или опросника для определения типа личности исходя из соотношения тех или иных черт;
	\item \textbf{К4} --- композиция (обобщение) других диспозициальных подходов;
	\item \textbf{К5} --- кросс-культурная устойчивость.
\end{itemize}

\begin{table}[h]
	\centering
	\begin{threeparttable}
		\caption{Сравнение подходов к описанию персональных черт}
		\label{tbl:exist_sol_transposed}
		\begin{tabularx}{\textwidth}{|X|X|p{3cm}|X|X|p{4cm}|}
			\hline
			\textbf{Подход} & \textbf{К1} & \textbf{К2} & \textbf{К3} & \textbf{К4} & \textbf{К5} \\
			\hline
			Олпорт & -- & Высокая & -- & -- & Неустойчива~\cite{kelland2025allport} \\
			\hline
			Айзенк & + & Высокая & + & -- & Устойчива~\cite{barrett1998epq, berry2002traits} \\
			\hline
			Кеттел & -- & Высокая & + & -- & Неустойчива~\cite{abrahams1996crosscultural} \\
			\hline
			БПЧ & + & Высокая & + & + & Устойчива~\cite{ahmad2021framework, berry2002traits} \\
			\hline
			MBTI & -- & Ограничена 16 типами & + & -- & Неустойчива~\cite{yoo_robinson_mbti_multicultural} \\
			\hline
		\end{tabularx}
	\end{threeparttable}
\end{table}

Ввиду наличия у некоторых описанных выше подходов конечных наборов факторов, которые можно рассматривать с различной степенью выраженности, можно сказать, что все рассмотренные ПОПЧ, за исключением MBTI и теории Олпорта, позволяют описывать произвольное число личностей на основе варьирования значений конкретных параметров.
Так, конкретная черта может быть задана как бинарная, тернарная, на основе иной шкалы или на основе соотношения вклада черт в характер~\cite{cattell_questionnaire, eysenck_questionnaire, big_five_questionairre}.
Для дальнейшей работы наиболее предпочтительной является БПЧ, поскольку она обладает высокой вариативностью, имеет развитую систему опросников, обобщает другие диспозициальные подходы и используется в исследованиях ДС и социальных роботов.

\section{Учёт персональных черт}

\subsection{Фреймворк для создания синтетических личностей}

В работе~\cite{Recchiuto2023SoftwareFramework} предлагается программная архитектура для создания разнообразных личностей робота (рисунок~\ref{img:CEA-taxonomy-arch}).

\includesvgimage
{CEA-taxonomy-arch}
{f}
{h}
{0.75\textwidth}
{Программная архитектура для создания синтетических личностей}

Основными компонентами в этой архитектуре являются генератор личностей, диспетчер действий и планировщик.
Для создания различных личностей авторами была использована таксономия трёх черт из <<Большой пятёрки>>: экстраверсия, добросовестность и доброжелательность.
Создаваемая личность рассматривается как вектор в трёхмерном пространстве:
\begin{equation}
	\label{equ:personality}
	Personality = W_e \cdot E + W_a \cdot A + W_c \cdot C, 
\end{equation}
где

\begin{itemize}
	\item $E$ --- базисный вектор, соответствующий экстраверсии;
	\item $A$ --- базисный вектор, соответствующий доброжелательности;
	\item $C$ --- базисный вектор, соответствующий сознательности;
	\item $W_e$ --- вес черты экстраверсия;
	\item $W_a$ --- вес черты доброжелательность;
	\item $W_c$ --- вес черты сознательность.
\end{itemize}


%\begin{itemize}
%	\item
%\end{itemize}

Диспетчер действий использует распределение весов и для каждого действия случайным образом выбирает черту для реализации.
Генератор личностей на основе $BERT$~\cite{devlin2019bert} принимает выбранную черту и действие, затем формирует параметры поведения: стиль речи, высоту тона, громкость, скорость речи, взгляд, движения головы, амплитуду и скорость жестов, скорость перемещения, дистанцию.
Таким образом, черты личности переводятся в вербальные, паравербальные и невербальные параметры реакции.

\subsection{PERSONAGE}

PERSONAGE (PERSONAlity GEnerator)~\cite{Mairesse2007Personage} --- параметрический генератор для порождения текстов с учётом одного из факторов <<Большой пятёрки>> --- экстраверсии.
В отличие от описанного выше фреймворка, PERSONAGE формирует только текстовую сторону реакции.
За основу был взят SPaRKy generator~\cite{Walker2002SPaRKy}, использовавшийся для генерации рекомендаций ресторанов Нью-Йорка.

Авторы адаптировали психолингвистические исследования экстраверсии к генератору в виде 29 параметров: 10 параметров для модуля планирования содержания и 19 параметров для модуля планирование предложений~\cite{Mairesse2007Personage}.
Эти параметры управляют объёмом сообщения, отбором позитивной и негативной информации, порядком представления фактов, сложностью синтаксических конструкций, количеством и частотой использования смягчающих слов и выражений.
Эмпирическая оценка показала, что тексты, сгенерированные с различными настройками параметров, статистически различаются по воспринимаемой экстраверсии и естественности речи.

\subsection{Генератор эмоционального состояния и выражений лица на основе входных фраз}

В работе~\cite{Chae2025AffectiveState} представлена методология, позволяющая генерировать эмоции и выражения лица на основе входной фразы самого робота и его личности.
В основе методологии лежит представление аффективных состояний в виде двумерного пространства, две оси которого составляют удовольствие (pleasure) и возбуждение (arousal)~\cite{russell1980circumplex}.
Для формирования личности робота были использованы черты <<Большой пятёрки>>, отображаемые в это пространство при помощи следующих уравнений~\cite{mehrabian1996pad}:
\begin{align}
	\alpha_p &= 0.21 \cdot E + 0.59 \cdot A - 0.19 \cdot N,\\
	\beta_p  &= 0.15 \cdot O + 0.3 \cdot A + 0.57 \cdot N,
\end{align}
где

\begin{itemize}
	\item $E$ --- экстраверсия;
	\item $O$ --- открытость опыту;
	\item $A$ --- доброжелательность;
	\item $N$ --- нейротизм;
	\item $\alpha_p$ --- значение удовольствия;
	\item $\beta_p$ --- значение возбуждения.
\end{itemize}

Далее вычисленные параметры личности используются совместно с аффективным анализом входной фразы для определения целевого эмоционального состояния робота и последующей генерации плавного эмоционального перехода и мимических выражений.

\subsection{Роботизированная голова KMC-EXPR}

В работе~\cite{park2012lawofattraction} исследовался вопрос о том, проявляется ли <<закон притяжения>> (similarity-attraction) в контексте взаимодействия человека и робота.
В исследовании использовалась роботизированная голова KMC-EXPR, содержащая подвижные элементы рта, губ и глаз, что позволяет роботу выражать широкий спектр мимических конфигураций.

Авторами были реализованы две противоположные выраженности черты экстраверсии из модели <<Большой пятёрки>>.
В экстравертном режиме робот проявлял активность и заинтересованность в общении: выполнял более широкие и быстрые движения глаз и рта, поддерживал частый зрительный контакт, оперативно реагировал на действия человека.
В интровертном режиме робот демонстрировал более пассивный и сдержанный стиль поведения: избегал прямого зрительного контакта, выполнял уменьшенные по амплитуде и более медленные движения рта и глаз.
Следовательно, личность робота была задана через невербальные параметры поведения.

\subsection{Робот-помощник для людей, переживших инсульт}

В работе~\cite{tapus2008personalitymatching} описывается исследование соответствия личности пользователя и робота-терапевта, созданного для помощи людям, пережившим инсульт, в выполнении реабилитационных упражнений.
Робот реализует нетактильный подход: он не прикасается к пациенту, а оказывает поддержку через социальное поведение, голосовые подсказки, регулирование дистанции и стимулирующее общение.

Для формирования личности робота авторы использовали модель личности Айзенка, из которой была выделена ключевая черта --- экстраверсия.
В контексте робота-терапевта эта черта выражалась через дистанцию взаимодействия, скорость и амплитуду движений, а также вербальную и паравербальную коммуникацию.
Интровертный режим соответствовал поддерживающему и мягкому стилю общения, экстравертный --- более активному и мотивирующему стилю.

\subsection{Модель на основе нечёткой логики}

В работе~\cite{xue2017fuzzy} предложена многоагентная модель пешеходного поведения, в которой персональные черты используются для получения неоднородного поведения агентов.
Несмотря на то, что работа не посвящена диалоговым роботам, она представляет интерес как способ переноса черт личности в параметры поведения с помощью нечёткой логики.

Для описания личности используется модель OCEAN, то есть те же пять факторов БПЧ: открытость опыту, добросовестность, экстраверсия, доброжелательность и нейротизм.
Каждый фактор задаётся числом в диапазоне $[-100, 100]$, а нейтральная личность соответствует вектору $(0, 0, 0, 0, 0)$.
Далее вводится понятие параметры принятия решений: максимальная дистанция учёта соседей, число учитываемых соседей, горизонт планирования, радиус личного пространства и предпочитаемая скорость.
Нечёткая система вывода устанавливает соответствие между факторами OCEAN и этими параметрами.
Например, экстраверсия и доброжелательность влияют на радиус личного пространства, а экстраверсия и нейротизм --- на предпочитаемую скорость.

Данный подход показывает, что черты личности не обязательно должны напрямую сопоставляться с готовыми реакциями.
Они могут задавать промежуточные управляющие параметры, которые затем используются исполнительной моделью.

\clearpage

\section{Сравнение способов учёта персональных черт}

В таблице~\ref{tbl:exist_sol_transposed2} представлены рассмотренные решения.
В качестве критериев сравнения способов учёта персональных черт были взяты следующие:

\begin{itemize}
	\item \textbf{K1} --- выбранный ПОПЧ;
	\item \textbf{K2} --- выбранные черты подхода;
	\item \textbf{K3} --- каналы коммуникации или поведенческие параметры.
\end{itemize}

\begin{table}[h]
	\centering
	\begin{threeparttable}
		\caption{Сравнение рассмотренных решений}
		\label{tbl:exist_sol_transposed2}
		\begin{tabularx}{\textwidth}{|X|X|X|X|}
			\hline
			\textbf{Решение} & \textbf{К1} & \textbf{К2} & \textbf{К3} \\
			\hline
			Генератор синтетических личностей~\cite{Recchiuto2023SoftwareFramework} &
			<<Большая пятёрка черт>> &
			Экстраверсия, добросовестность, доброжелательность &
			вербальный, паравербальный, невербальный \\
			\hline
			PERSONAGE~\cite{Mairesse2007Personage} &
			<<Большая пятёрка черт>> &
			Экстраверсия &
			вербальный \\
			\hline
			Генератор эмоционального состояния~\cite{Chae2025AffectiveState} &
			<<Большая пятёрка черт>> &
			Экстраверсия, доброжелательность, нейротизм, открытость опыту &
			невербальный \\
			\hline
			Роботизированная голова KMC-EXPR~\cite{park2012lawofattraction} &
			<<Большая пятёрка черт>> &
			Экстраверсия &
			невербальный \\
			\hline
			Робот-помощник для людей, переживших инсульт~\cite{tapus2008personalitymatching} &
			Модель Айзенка &
			Экстраверсия &
			вербальный, паравербальный, невербальный \\
			\hline
			Модель на основе нечёткой логики~\cite{xue2017fuzzy} &
			<<Большая пятёрка черт>> &
			Все черты &
			пространственные и двигательные параметры \\
			\hline
		\end{tabularx}
	\end{threeparttable}
\end{table}

По данным таблицы~\ref{tbl:exist_sol_transposed2}, экстраверсия используется почти всеми рассмотренными решениями.
Базисным подходом большинства исследований является <<Большая пятёрка черт>>, при этом это единственная типология из рассмотренных в таблице~\ref{tbl:exist_sol_transposed}, обобщающая другие подходы к описанию персональных черт и обладающая кросс-культурной устойчивостью.
Здесь 5 основных черт с возможностью подразделения каждой на 6 граней, что при возможности варьирования вклада каждой грани как действительной величины позволяет описать множество характеров большой мощности.

Для модификации метода ведения диалога роботом Ф-2 целесообразно использовать БПЧ как модель описания личности, а способ учёта персональных черт строить как отображение значений черт в параметры мультимодальной реакции.

\phantomsection
\section*{Выводы}
\addcontentsline{toc}{section}{Выводы}

По результатам анализа в качестве базовой модели описания личности выбрана <<Большая пятёрка черт>>, поскольку она обладает высокой вариативностью, поддерживается опросниками, используется в работах по робототехнике и диалоговым агентам, а также имеет кросс-культурную устойчивость.

Формализованная постановка задачи в нотации IDEF0 представлена на рисунке~\ref{img:fpz_analysis}. % че то про реструктуризацию добавить надо

\includesvgimage
{fpz_analysis}
{f}
{h}
{1\textwidth}
{Формализованная постановка задачи в нотации IDEF0}
